\section{Ambiguitas Grammar dan Dangling Else}
\label{sec:ch09-ambiguitas}

\subsection{Definisi Ambiguitas}

\begin{definisi}[Tata bahasa ambigu]
Tata bahasa \(G\) disebut ambigu jika terdapat string
\(w \in L(G)\) yang memiliki dua pohon urai berbeda.
\end{definisi}

Ambiguitas berarti parser tidak dapat menentukan struktur yang
dimaksud tanpa aturan tambahan. Pada bahasa pemrograman, ambiguitas
harus dihilangkan atau diselesaikan dengan aturan eksplisit.

\begin{contoh}[Grammar ambigu]
Tata bahasa
\[
E \to E + E \mid E * E \mid \code{id}
\]
ambigu. String \(\code{id} + \code{id} * \code{id}\) memiliki dua
pohon urai: satu mengutamakan penjumlahan, satu mengutamakan
perkalian. Tanpa aturan prioritas, makna ekspresi tidak jelas.
\end{contoh}

\subsection{Solusi Ambiguitas}

Ada tiga cara umum menangani ambiguitas:
\begin{enumerate}[label=\arabic*.,leftmargin=*]
  \item \textbf{restrukturisasi grammar}: membuat tingkat prioritas
        terpisah;
  \item \textbf{aturan precedence dan associativity}: dipakai oleh
        parser generator;
  \item \textbf{aturan semantik}: memilih pohon berdasarkan konteks,
        misalnya else terdekat.
\end{enumerate}

Restrukturisasi grammar adalah solusi paling bersih karena
menghilangkan ambiguitas pada tingkat sintaksis. Grammar aritmetika
pada Bab VIII adalah contoh restrukturisasi: \(E\), \(T\), dan
\(F\) memisahkan tingkat prioritas.

\subsection{Dangling Else}

Ambiguitas klasik lain adalah \eng{dangling else}:
\[
\text{stmt} \to \text{if}\ E\ \text{then}\ \text{stmt}
\mid \text{if}\ E\ \text{then}\ \text{stmt}\ \text{else}\
\text{stmt} \mid \text{other}.
\]
String
\[
\text{if}\ E_1\ \text{then}\ \text{if}\ E_2\ \text{then}\ S_1\
\text{else}\ S_2
\]
dapat diurai dengan dua cara: \code{else} milik \code{if} luar atau
\code{if} dalam.

\begin{figure}[H]
\centering
\begin{minipage}{0.46\textwidth}
\centering
\begin{forest}
parsetree
[if
  [\(E_1\)]
  [then
    [if
      [\(E_2\)]
      [then [\(S_1\)]]
      [else [\(S_2\)]]
    ]
  ]
]
\end{forest}
\end{minipage}\hfill
\begin{minipage}{0.46\textwidth}
\centering
\begin{forest}
parsetree
[if
  [\(E_1\)]
  [then
    [if
      [\(E_2\)]
      [then [\(S_1\)]]
    ]
  ]
  [else [\(S_2\)]]
]
\end{forest}
\end{minipage}
\caption{Dua pohon urai untuk dangling else}
\label{fig:ch09-dangling}
\end{figure}

Aturan yang dipakai hampir semua bahasa adalah \textbf{else terdekat}:
\code{else} dipasangkan dengan \code{if} terdekat yang belum memiliki
\code{else}. Aturan ini dapat dinyatakan dengan grammar
\code{matched}/\code{unmatched} atau diselesaikan oleh parser dengan
prioritas shift.

\begin{warningbox}
Ambiguitas bukan hanya masalah teoretis. Parser yang memilih pohon
secara acak dapat menghasilkan program yang berperilaku berbeda dari
yang diharapkan. Selalu tetapkan aturan resolusi dan uji dengan kasus
dangling else.
\end{warningbox}
